\documentclass[italian, oneside]{sapthesis}
\usepackage[utf8]{inputenc}
\usepackage[T1]{fontenc}
\usepackage{babel}
\usepackage[hidelinks]{hyperref}
\usepackage{amsmath, amssymb}
\usepackage{graphicx}
\usepackage{booktabs}
\usepackage{microtype}
\usepackage{caption}
\usepackage{subcaption}

% --- Dati del candidato e della tesi ---
\title{Orchestratore ILP centralizzato per il routing di task su reti satellitari}
\author{Claudio Bagini}
\IDnumber{2045337}
\course{Informatica}
\courseorganizer{Facoltà di Ingegneria dell'Informazione, Informatica e Statistica}
\submitdate{2025/2026}
\copyyear{2026}
\advisor{Dott. Emiliano Casalicchio}
\coadvisor{}
\authoremail{bagini.2045337@studenti.uniroma1.it}
\versiondate{\today}
\thesistype{Tesi di Laurea Triennale}

%\examdate{Data della discussione}
%\examiner{Presidente della commissione}
%\examiner[Relatore]{Dott. Emiliano Casalicchio}
%\examiner{Membro della commissione}

\begin{document}

\frontmatter
\maketitle

\dedication{Dedica}

\begin{abstract}
L'evoluzione delle megacostellazioni di satelliti in orbita bassa (LEO, \emph{Low Earth Orbit}) ha aperto nuove prospettive per il \emph{Satellite Computing}, consentendo l'estensione delle capacità computazionali dell'edge computing direttamente nello spazio (\emph{Orbital Edge Computing}). Tuttavia, l'elevata dinamicità della topologia di rete, i vincoli energetici stringenti dei nodi satellitari e le rigorose scadenze temporali (\emph{deadline}) delle applicazioni \emph{real-time} rendono l'orchestrazione e l'offloading dei task sfide di primaria importanza.

Questo lavoro di tesi presenta la progettazione, l'implementazione e l'analisi prestazionale di un orchestratore centralizzato basato sulla Programmazione Lineare Intera (ILP). Per mitigare i costi computazionali del solutore e prevenire fenomeni di \emph{starvation}, il sistema integra un meccanismo di \emph{batching} dinamico regolato dalle metriche di \emph{Batch Size} ($BS$) e \emph{Batch Timeout} ($BT$), insieme a una formulazione coordinata dei vincoli di temporizzazione che compensa i tempi di permanenza in coda prima della risoluzione.

Attraverso simulazioni a eventi discreti sviluppate nel framework SimPy, l'architettura è stata valutata sotto stress al variare dei tassi di arrivo delle richieste ($\lambda$), confrontando le strategie di ottimizzazione \emph{Mapping Time} e \emph{Mapping Energy}. L'analisi dettagliata ha evidenziato che i punti di collasso sotto carico intenso ($\lambda > 6 \text{ req/s}$) sono prevalentemente causati dalla scadenza dei tempi di attesa (\emph{Deadline Exceeded}) piuttosto che dalla saturazione delle risorse fisiche di calcolo. Inoltre, gli esperimenti di sensibilità hanno dimostrato l'efficacia del rilassamento delle deadline ($\Delta D$) nel recuperare tassi di completamento elevati.

Infine, l'orchestratore centralizzato è stato sottoposto a un benchmark comparativo con le soluzioni distribuite ed euristiche dello stato dell'arte (\emph{OrbitAware}, \emph{DTS-base}, \emph{DTS-APopt} e ILP distribuito). I risultati quantificano il \emph{trade-off} fondamentale tra l'ottimalità globale teorica dell'allocazione centralizzata e l'impatto della latenza di accodamento rispetto alla rapidità di risposta dei modelli distribuiti.
\end{abstract}

\begin{acknowledgments}
Ringraziamenti
\end{acknowledgments}

\tableofcontents

\mainmatter

% --- INCLUSIONE DEI CAPITOLI SEPARATI ---
\chapter{Introduzione}
\label{chap:introduzione}

\section{Contesto e Motivazioni}


\section{Formulazione del Problema}


\section{Obiettivo della Tesi}


\section{Organizzazione del Testo}

\chapter{Modello di Sistema e Stato dell'Arte}
\label{chap:stato_arte}

\section{Modello di Reti Satellitari e Workload}
\subsection{Architettura di Rete e Collegamenti ISL}
\subsection{Caratteristiche del Carico di Lavoro}

\section{Tassonomia delle Soluzioni di Orchestrazione}
\subsection{Approcci Centralizzati vs Distribuiti}
\subsection{Punti di Forza e Limiti delle Architetture Attuali}

\section{Algoritmi di Riferimento dello Stato dell'Arte}
\subsection{Euristiche Distribuite: OrbitAware, DTS-base, DTS-APopt}
\subsection{Modello ILP Gerarchico e Distribuito di Riferimento}
\chapter{Architettura dell'Orchestratore Centralizzato ILP}
\label{chap:architettura_ilp}

\section{Formulazione del Modello ILP}
\subsection{Funzioni Obiettivo: Mapping Time e Mapping Energy}
\subsection{Vincoli di Risorsa, Rete ed Energia}

\section{Meccanismo di Batching Integrato}
\subsection{Gestione tramite Batch Size e Batch Timeout}
\subsection{Mitigazione dei Fenomeni di Starvation}

\section{Modellazione Dinamica delle Deadline}
\subsection{Aggiustamento del Tempo di Attesa in Coda}
\chapter{Analisi delle Prestazioni e Sensibilità dell'Orchestratore ILP}
\label{chap:analisi_prestazioni}

\section{Setup Sperimentale}
\subsection{Ambiente di Simulazione a Eventi Discreti (SimPy)}
\subsection{Parametri di Rete e Profilo Hardware del Sistema}

\section{Impatto dei Parametri di Batching e Timeout}

\section{Analisi Dettagliata delle Prestazioni}
\subsection{Tasso di Completamento (Success Rate) e Punti di Collasso}
\subsection{Distribuzione e Diagnosi delle Cause di Rigetto}
\subsection{Scomposizione dei Tempi di Risposta e di Sistema}

\section{Sensibilità al Rilassamento delle Deadline}
\chapter{Valutazione Comparativa: Orchestratore Centralizzato vs Soluzioni Distribuiti}
\label{chap:benchmark}

\section{Setup del Benchmark Comparativo}

\section{Confronto delle Prestazioni}
\subsection{Success Rate e Latenza di Risposta}
\subsection{Grado di Bilanciamento del Carico (Load Balancing Degree)}
\subsection{Efficienza Energetica Complessiva}

\section{Discussione dei Trade-off}
\subsection{Latenza di Aggregazione vs Ottimalità Globale}
\chapter{Conclusioni e Sviluppi Futuri}
\label{chap:conclusioni}

\section{Sintesi dei Risultati}

\section{Limiti dell'Approccio Centralizzato}

\section{Sviluppi Futuri e Architetture Ibride}

\backmatter
\cleardoublepage
\phantomsection
\addcontentsline{toc}{chapter}{Bibliografia}
\nocite{*}
\bibliographystyle{sapthesis}
%\bibliography{bibliografia}

\end{document}